\documentclass[letterpaper,11pt]{article}

\usepackage[utf8]{inputenc}
\usepackage[T1]{fontenc}
\usepackage[spanish,es-nodecimaldot]{babel}
\usepackage{amsmath,amssymb}
\usepackage{geometry}
\usepackage{graphicx}
\usepackage{xcolor}
\usepackage{fancyhdr}
\usepackage{booktabs}
\usepackage{tabularx}
\usepackage{array}
\usepackage{tcolorbox}
\usepackage{enumitem}
\usepackage{lastpage}
\tcbuselibrary{breakable,skins}

\geometry{letterpaper,left=1.8cm,right=1.8cm,top=2.5cm,bottom=1.8cm,headheight=1.6cm}

\definecolor{azul}{HTML}{174A70}
\definecolor{celeste}{HTML}{EDF5FA}
\definecolor{gris}{HTML}{5F6872}
\definecolor{verde}{HTML}{24724A}

\newcommand{\logousach}{../../template/images/logo_usach.png}
\newcommand{\logoup}{../../template/images/logo_up.png}
\newcommand{\R}{\operatorname{Re}}
\newcommand{\Ima}{\operatorname{Im}}

\pagestyle{fancy}
\fancyhf{}
\fancyhead[L]{\includegraphics[height=1.05cm]{\logousach}}
\fancyhead[R]{\begin{minipage}[b]{0.62\textwidth}\raggedleft
  \small\textbf{\color{azul}Álgebra I}\\
  \footnotesize Control 1 - Forma C
\end{minipage}}
\fancyfoot[L]{\footnotesize Documento de apoyo académico}
\fancyfoot[R]{\footnotesize Página \thepage\ de \pageref{LastPage}}
\renewcommand{\headrulewidth}{0.5pt}
\renewcommand{\footrulewidth}{0.4pt}

\newtcolorbox{enunciado}[1]{
  enhanced,breakable,colback=celeste,colframe=azul,
  colbacktitle=azul,coltitle=white,fonttitle=\bfseries,
  title={#1},arc=3pt,boxrule=0.8pt,left=8pt,right=8pt,top=4pt,bottom=4pt
}
\newtcolorbox{desarrollo}[1]{
  enhanced,breakable,colback=white,colframe=azul!65,
  colbacktitle=celeste,coltitle=azul,fonttitle=\bfseries,
  title={#1},arc=3pt,boxrule=0.7pt,left=9pt,right=9pt,top=4pt,bottom=4pt
}
\newtcolorbox{respuesta}{
  colback=verde!7,colframe=verde,arc=3pt,boxrule=0.9pt,
  left=8pt,right=8pt,top=4pt,bottom=4pt
}

\setlength{\parindent}{0pt}
\setlength{\parskip}{5pt}
\renewcommand{\arraystretch}{1.18}

\begin{document}
\thispagestyle{empty}

\begin{minipage}[c]{0.18\textwidth}
  \includegraphics[width=2.35cm]{\logousach}
\end{minipage}
\begin{minipage}[c]{0.64\textwidth}
  \centering
  {\large\bfseries\color{azul} Álgebra I}\\
  {\small Control 1 - Forma C}
\end{minipage}
\begin{minipage}[c]{0.17\textwidth}
  \raggedleft\includegraphics[width=2.0cm]{\logoup}
\end{minipage}

\vspace{1.8cm}

\begin{center}
  {\color{azul}\rule{0.82\textwidth}{1.2pt}}\\[0.5cm]
  {\Huge\bfseries\color{azul} PAUTA DE CORRECCIÓN}\\[0.35cm]
  {\LARGE\bfseries Control 1}\\[0.2cm]
  {\Large Forma C}\\[0.45cm]
  {\color{azul}\rule{0.82\textwidth}{1.2pt}}
\end{center}

\vspace{1.15cm}

\begin{tcolorbox}[colback=celeste,colframe=azul,arc=4pt,boxrule=0.9pt]
\centering
\begin{tabular}{rl}
\textbf{Contenido:} & Números complejos\\
\textbf{Duración del control:} & 50 minutos
\end{tabular}
\end{tcolorbox}

\vfill
\begin{center}
  \small\color{gris}Pauta desarrollada para los dos ejercicios del Control 1, Forma C.
\end{center}

\newpage
\small
\setlength{\parskip}{2pt}
\setlength{\abovedisplayskip}{4pt}
\setlength{\belowdisplayskip}{4pt}
\setlength{\abovedisplayshortskip}{3pt}
\setlength{\belowdisplayshortskip}{3pt}

\begin{enunciado}{Ejercicio 1}
Sean \(z,w\in\mathbb{C}\) dados por
\[
 z=\frac{1-7i}{2+i},
 \qquad
 w=\frac{6+2i}{-i+2}.
\]
Determine el valor de
\[
 \R\!\left(3\overline z-2\overline w\right)+\Ima\!\left(z^3\right).
\]
\end{enunciado}

\begin{desarrollo}{1. Expresar \(z\) y \(w\) en forma binómica}
Racionalizamos cada denominador usando su conjugado:
\[
\begin{aligned}
z
&=\frac{1-7i}{2+i}\cdot\frac{2-i}{2-i}
 =\frac{(1-7i)(2-i)}{(2+i)(2-i)}\\
&=\frac{2-i-14i+7i^2}{2^2+1^2}
 =\frac{-5-15i}{5}
 =-1-3i.
\end{aligned}
\]
\[
\begin{aligned}
w
&=\frac{6+2i}{2-i}\cdot\frac{2+i}{2+i}
 =\frac{(6+2i)(2+i)}{(2-i)(2+i)}\\
&=\frac{12+6i+4i+2i^2}{2^2+1^2}
 =\frac{10+10i}{5}
 =2+2i.
\end{aligned}
\]
\end{desarrollo}

\begin{desarrollo}{2. Calcular la parte real solicitada}
Como
\[
\overline z=-1+3i,
\qquad
\overline w=2-2i,
\]
se obtiene
\[
\begin{aligned}
3\overline z-2\overline w
&=3(-1+3i)-2(2-2i)\\
&=(-3+9i)-(4-4i)\\
&=-7+13i.
\end{aligned}
\]
Por lo tanto,
\[
\R\!\left(3\overline z-2\overline w\right)=-7.
\]
\end{desarrollo}

\begin{desarrollo}{3. Calcular la parte imaginaria de \(z^3\)}
Primero,
\[
z^2=(-1-3i)^2=1+6i+9i^2=-8+6i.
\]
Luego,
\[
\begin{aligned}
z^3
&=(-8+6i)(-1-3i)\\
&=8+24i-6i-18i^2\\
&=26+18i.
\end{aligned}
\]
En consecuencia,
\[
\Ima(z^3)=18.
\]
\end{desarrollo}

\begin{respuesta}
\textbf{Respuesta final.}
\[
\boxed{\R\!\left(3\overline z-2\overline w\right)+\Ima(z^3)
=-7+18=11}
\]
\end{respuesta}

\newpage

\begin{enunciado}{Ejercicio 2}
Sea \(z\in\mathbb{C}\) tal que
\[
z^2-2+4i\sqrt2=0.
\]
Determine los valores de \(\R(z)\) e \(\Ima(z)\).
\end{enunciado}

\begin{desarrollo}{1. Aislar el cuadrado y plantear la forma binómica}
De la ecuación,
\[
z^2=2-4i\sqrt2.
\]
Sea \(z=a+bi\), con \(a,b\in\mathbb{R}\). Entonces
\[
z^2=(a+bi)^2=(a^2-b^2)+2ab\,i.
\]
Igualando partes real e imaginaria:
\[
\boxed{a^2-b^2=2},
\qquad
\boxed{2ab=-4\sqrt2},
\]
es decir, \(ab=-2\sqrt2\).
\end{desarrollo}

\begin{desarrollo}{2. Determinar \(a\) y \(b\)}
El módulo de \(z^2\) satisface
\[
|z^2|=\left|2-4i\sqrt2\right|
=\sqrt{2^2+\left(-4\sqrt2\right)^2}
=\sqrt{4+32}=6.
\]
Pero \(|z^2|=|z|^2=a^2+b^2\); por tanto,
\[
a^2+b^2=6.
\]
Sumando esta igualdad con \(a^2-b^2=2\):
\[
2a^2=8
\quad\Longrightarrow\quad
a^2=4
\quad\Longrightarrow\quad
a=\pm2.
\]
Usamos \(ab=-2\sqrt2\):
\[
\begin{array}{c|c|c}
a & b=-\dfrac{2\sqrt2}{a} & z=a+bi\\ \midrule
2 & -\sqrt2 & 2-\sqrt2\,i\\
-2 & \sqrt2 & -2+\sqrt2\,i
\end{array}
\]
\end{desarrollo}

\begin{desarrollo}{3. Comprobación algebraica}
\[
\begin{aligned}
(2-\sqrt2\,i)^2
&=4-4\sqrt2\,i+2i^2
=2-4\sqrt2\,i,\\
(-2+\sqrt2\,i)^2
&=4-4\sqrt2\,i+2i^2
=2-4\sqrt2\,i.
\end{aligned}
\]
Ambos valores satisfacen la ecuación original.
\end{desarrollo}

\begin{respuesta}
\textbf{Respuesta final.} Existen dos soluciones:
\[
\boxed{
\begin{array}{c|cc}
z & \R(z) & \Ima(z)\\ \midrule
2-\sqrt2\,i & 2 & -\sqrt2\\
-2+\sqrt2\,i & -2 & \sqrt2
\end{array}}
\]
\end{respuesta}

\end{document}
